%% ma: L12-B08-0049
%% lop: 12
%% chuong: 2
%% bai: 8
%% dang: 12.8.33
%% loai: DS
%% muc: [NB, TH, TH, VD]
%% dap_an: "ĐSSS"
%% nguon: "LaTeX giáo viên gửi 21/10/2026 (M05, Vectơ và hệ trục tọa độ trong không gian, 107 câu), Phần II Câu 1"
%% kien_thuc: "Bài 8 (tọa độ vectơ, tích vô hướng, hai vectơ cùng phương, vuông góc)"
%% trang_thai: cho-duyet
%% ly_do: "Dạng mới đề xuất 12.8.33"
%% ghi_chu: "Ý d) đề gốc cho H(47/29; 13/29; 4/29) không phải chân đường cao (đúng là H(-3/7; -9/14; 39/14)); ý d) vẫn xác định được là Sai"
\begin{ex}
Trong không gian $Oxyz$, cho ba điểm $A(1;0;-2)$, $B(2;3;4)$ và $C(-4;-6;1)$.
\choiceTF
{\True Tọa độ vectơ $\overrightarrow{AB}$ là $(1;3;6)$.}
{Tích vô hướng $\overrightarrow{AB}\cdot\overrightarrow{AC}=0$.}
{Tam giác $ABC$ là tam giác vuông tại $A$.}
{Chân đường cao hạ từ đỉnh $A$ của tam giác $ABC$ có tọa độ là $H\left(\dfrac{47}{29};\dfrac{13}{29};\dfrac4{29}\right)$.}
\loigiai{a) $\overrightarrow{AB}=(1;3;6)$. Đúng.

b) $\overrightarrow{AC}=(-5;-6;3)$, $\overrightarrow{AB}\cdot\overrightarrow{AC}=-5-18+18=-5\ne0$. Sai.

c) $\overrightarrow{AB}\cdot\overrightarrow{AC}\ne0$ nên tam giác không vuông tại $A$. Sai.

d) $\overrightarrow{BC}=(-6;-9;-3)$ cùng phương với $\vec{u}=(2;3;1)$. Gọi $H=B+t\vec{u}=(2+2t;3+3t;4+t)$. $\overrightarrow{AH}\cdot\vec{u}=0\Leftrightarrow17+14t=0\Leftrightarrow t=-\dfrac{17}{14}$, suy ra $H\left(-\dfrac37;-\dfrac9{14};\dfrac{39}{14}\right)$. Sai.}
\end{ex}
